\section{模型的评价与推广}

% 参考优秀论文 A196 的末章组织：评价必须对应前文已经得到的模型、结果和验证证据，
% 不是统一套用“优点很多、推广性强”等空泛表述。

\subsection{模型的优点}
% 通常约 3 项有据优势，兼顾模型结构与求解性质；数量不作为自动配额。
% 正式实例须据真实模型和结果替换占位，不得编造效率、精度和鲁棒性。
\begin{enumerate}
    \item ……。
\end{enumerate}

\subsection{模型的不足与改进方向}
% 通常 1–2 项限制，说明受影响的结果、边界及对应改进方向，不强制凑条目。
\begin{enumerate}
    \item ……。
\end{enumerate}

% \subsection{模型的推广}
% 选写：确有可迁移对象时，用短段说明可继承结构、需重新标定的参数/边界。

% 中文国赛默认不单设“结论”一级章；各问题的最终回答应已经在对应“模型求解结果与分析”中闭环。
